Dos pequeños depósitos forman una cascada en un ensayo simulado: el caudal que sale del primero entra íntegramente en el segundo, sin retorno ni retardo. Al primero entra \(1\,\mathrm{L/min}\). El caudal de salida de cada depósito es su volumen dividido por un tiempo característico: \(2\,\mathrm{min}\) para el primero y \(1\,\mathrm{min}\) para el segundo. No hay fugas ni rebose en el intervalo estudiado. Inicialmente contienen \(4\,\mathrm L\) y \(1\,\mathrm L\), respectivamente.

Se quiere predecir la evolución mediante una ecuación diferencial de segundo orden cuya única variable dependiente sea el volumen del segundo depósito. Elija las variables, construya el modelo y determine sus condiciones iniciales a partir de la información física. Obtenga la evolución del segundo volumen y reconstruya la del primero para comprobar ambos balances.
